% LỜI TỰA --- Quyển XIX
\chapter*{Lời Tựa}
\addcontentsline{toc}{chapter}{Lời Tựa}
\markboth{Lời Tựa}{Lời Tựa}

\begin{truyen}{Lời Tựa}{100 Sentences That Make Students Fall Silent}
\chuhoa{G}{iáo viên bắt đầu bộ sách này với một câu hỏi đơn giản: làm thế nào để mỗi buổi học có thể để lại thứ gì đó đáng nhớ hơn là bài thi?}
\textit{(A teacher began this series with a simple question: how can each class leave behind something more memorable than a test?)}

Có những câu thầy nói xong, cả lớp im lặng~--- không phải im vì sợ, mà im vì đang suy nghĩ.
\textit{(Some sentences the teacher says leave the whole class silent~--- not from fear, but because they are thinking.)}

Quyển mười chín gom 100 câu như thế~--- ngắn, sâu, khiến học sinh dừng lại và tự hỏi. Mỗi câu đi kèm một câu chuyện nhỏ và gợi ý để giữ lại. Mong rằng thầy cô có thêm câu để nói đúng lúc; học trò có thêm điều để im lặng mà nghĩ.
\textit{(Volume nineteen gathers 100 such sentences~--- short, deep, that make students pause and wonder. Each comes with a short story and a takeaway. May teachers have more sentences to say at the right moment; may students have more to ponder in silence.)}

Bộ sách <<Truyện Tích Cảm Hứng Song Ngữ>> gồm 28 quyển, mỗi quyển một chủ đề riêng. Quyển XIX này mang chủ đề: \textbf{những câu ngắn, sâu, khiến học sinh dừng lại và tự hỏi}.
\textit{(The series ``Inspiring Bilingual Tales'' contains 28 volumes, each with its own theme. Volume XIX focuses on: \textbf{những câu ngắn, sâu, khiến học sinh dừng lại và tự hỏi}.})
\end{truyen}

\begin{baihoc}
Im lặng sau khi nghe một câu đúng lúc là dấu hiệu của tư duy, không phải của thờ ơ.
\textit{(Silence after the right sentence is a sign of thinking, not of indifference.)}
\end{baihoc}

\begin{ghinhoanh}
\textbf{Short English to remember:}

Silence after a deep sentence means thinking has begun.

One hundred words: one hundred doors to reflection.
\end{ghinhoanh}

\begin{tuhocnhanh}
1. Câu nào trong 100 câu này khiến em im lặng lâu nhất? Vì sao? \textit{(Which of the 100 sentences made you silent the longest? Why?)}

2. Em đã từng tạo ra một khoảnh khắc im lặng suy nghĩ cho người khác chưa? \textit{(Have you ever created a moment of silent reflection for someone else?)}

3. Câu nào em muốn thầy cô nói với cả lớp vào đầu buổi học ngày mai? \textit{(Which sentence would you want your teacher to say to the class tomorrow morning?)}
\end{tuhocnhanh}

\begin{center}
\textit{\color{silver}``Im lặng đôi khi là câu trả lời hay nhất.''}
\end{center}

\begin{flushright}
\textbf{Tôi Nguyễn Văn Sang}\\
\textit{GV THPT Nguyễn Hữu Cảnh - TP HCM}
\end{flushright}

\ngancach
